\documentclass[journal]{IEEEtran}
\usepackage[utf8]{inputenc}
\usepackage[T1]{fontenc}
\usepackage{amsmath,amssymb,amsfonts}
\usepackage{graphicx}
\usepackage{booktabs}
\usepackage{tabularx}
\usepackage{multirow}
\usepackage{threeparttable}
\usepackage{siunitx}
\usepackage{microtype}
\usepackage{algorithm}
\usepackage{algpseudocode}
\usepackage{xcolor}
\usepackage[hidelinks]{hyperref}
\usepackage{cite}

\graphicspath{{figures/}}

\begin{document}

\title{Deterministic Software Architecture and Automated Governance Workflow for Public Administration Systems}

\author{First~Author,~\IEEEmembership{Member,~IEEE,}
        and~Second~Author,~\IEEEmembership{Senior~Member,~IEEE}%
\thanks{Manuscript received Month DD, YYYY; revised Month DD, YYYY.}%
\thanks{Authors are with the Department of Computer Science, University Name, City, Country.}}

\maketitle

\begin{abstract}
This template establishes a Q1 IEEE Transactions baseline with strict two-column typography, vector figure embedding, booktabs tables, microtype protrusion, and zero overfull horizontal box tolerance. Replace this abstract with a structured 150--250 word summary covering context, problem gap, proposed method, empirical evaluation, and key quantitative gains.
\end{abstract}

\begin{IEEEkeywords}
Software architecture, workflow automation, formal verification, empirical evaluation.
\end{IEEEkeywords}

\section{Introduction}
\IEEEPARstart{M}{odern} public administration platforms require deterministic workflow enforcement, verifiable audit trails, and reproducible technical documentation. Prior approaches frequently suffer from divergence between system implementation and published architectural specifications.

Our primary contributions are threefold:
\begin{enumerate}
  \item A configuration-driven state machine decoupling institutional SOP rules from core engine code.
  \item A hash-chained audit ledger ensuring tamper-evident governance verification.
  \item An empirical evaluation across representative workload profiles demonstrating low-latency compliance.
\end{enumerate}

\section{Methodology}
The end-to-end validation metric combines latency $L_i$ and compliance accuracy $A_i$ across $N$ evaluated service workflows:
\begin{equation}
  S = \frac{1}{N} \sum_{i=1}^{N} \left( \alpha A_i - \beta \frac{L_i}{L_{\max}} \right),
  \label{eq:score}
\end{equation}
where $\alpha, \beta \in (0, 1]$ weight correctness against bounded service-level response time $L_{\max}$.

\begin{table}[!t]
\centering
\caption{Quantitative Comparison Across Evaluated Baselines}
\label{tab:results}
\begin{tabular}{@{}lccc@{}}
\toprule
Method & Latency (ms) & Accuracy (\%) & Memory (MB) \\
\midrule
Baseline Monolith & 142.4 $\pm$ 8.1 & 91.2 & 412 \\
Rule-Engine v1    & 68.7 $\pm$ 4.3  & 96.4 & 248 \\
\textbf{Proposed} & \textbf{24.1 $\pm$ 1.9} & \textbf{99.4} & \textbf{164} \\
\bottomrule
\end{tabular}
\end{table}

\section{Experimental Results}
As reported in Table~\ref{tab:results}, the proposed architecture achieves a $5.9\times$ reduction in median latency while elevating compliance accuracy to $99.4\%$. All measurements adhere to Eq.~(\ref{eq:score}).

\section{Conclusion}
We presented a reproducible Q1 IEEE manuscript baseline ready for vector figures, BibTeX references, and automated visual quality verification.

\end{document}
